\documentclass{thesisclass}
% Based on thesisclass.cls of Timo Rohrberg, 2009
% ----------------------------------------------------------------
% Thesis - Main document
% ----------------------------------------------------------------


%% -------------------------------
%% |  Information for PDF file   |
%% -------------------------------
\hypersetup{
 pdfauthor={Not set},
 pdftitle={Not set},
 pdfsubject={Not set},
 pdfkeywords={Not set}
}


%% ---------------------------------
%% | Information about the thesis  |
%% ---------------------------------

\newcommand{\myname}{Name}
\newcommand{\mytitle}{Title of thesis\\ 
											Second title line}
\newcommand{\myinstitute}{Institute for Program Structures\\
													and Data Organization (IPD)}

\newcommand{\reviewerone}{?}
\newcommand{\reviewertwo}{?}
\newcommand{\advisor}{?}
\newcommand{\advisortwo}{?}

\newcommand{\timestart}{XX. Monat 20XX}
\newcommand{\timeend}{XX. Monat 20XX}
\newcommand{\submissiontime}{DD. MM. 20XX}


%% ---------------------------------
%% | ToDo Marker - only for draft! |
%% ---------------------------------
% Remove this section for final version!
\setlength{\marginparwidth}{20mm}

\newcommand{\margtodo}
{\marginpar{\textbf{\textcolor{red}{ToDo}}}{}}

\newcommand{\todo}[1]
{{\textbf{\textcolor{red}{(\margtodo{}#1)}}}{}}


%% --------------------------------
%% | Old Marker - only for draft! |
%% --------------------------------
% Remove this section for final version!
\newenvironment{deprecated}
{\begin{color}{gray}}
{\end{color}}


%% --------------------------------
%% | Settings for word separation |
%% --------------------------------
% Help for separation:
% In german package the following hints are additionally available:
% "- = Additional separation
% "| = Suppress ligation and possible separation (e.g. Schaf"|fell)
% "~ = Hyphenation without separation (e.g. bergauf und "~ab)
% "= = Hyphenation with separation before and after
% "" = Separation without a hyphenation (e.g. und/""oder)

% Describe separation hints here:
\hyphenation{
% Pro-to-koll-in-stan-zen
% Ma-na-ge-ment  Netz-werk-ele-men-ten
% Netz-werk Netz-werk-re-ser-vie-rung
% Netz-werk-adap-ter Fein-ju-stier-ung
% Da-ten-strom-spe-zi-fi-ka-tion Pa-ket-rumpf
% Kon-troll-in-stanz
}


%% ------------------------
%% |    Including files   |
%% ------------------------
% Only files listed here will be included!
% Userful command for partially translating the document (for bug-fixing e.g.)
\includeonly{%
titlepage,
declaration,
introduction,
content,
evaluation,
conclusion,
appendix
}


%%%%%%%%%%%%%%%%%%%%%%%%%%%%%%%%%
%% Here, main documents begins %%
%%%%%%%%%%%%%%%%%%%%%%%%%%%%%%%%%
\begin{document}

% Remove the following line for German text
\selectlanguage{english}

\frontmatter
\pagenumbering{roman}
%% titlepage.tex
%%
% coordinates for the bg shape on the titlepage
\selectlanguage{english}
\newcommand{\diameter}{20}
\newcommand{\xone}{-6}
\newcommand{\xtwo}{168}
\newcommand{\yone}{12}
\newcommand{\ytwo}{-255}

%\KOMAoptions{twoside=semi,DIV=last}

%~ \begin{titlepage}
%~ \selectlanguage{ngerman}
%~ % bg shape
%~ \begin{tikzpicture}[overlay]
%~ \draw[color=gray]  
 		 %~ (\xone mm, \yone mm)
  %~ -- (\xtwo mm, \yone mm)
 %~ arc (90:0:\diameter pt) 
  %~ -- (\xtwo mm + \diameter pt , \ytwo mm) 
	%~ -- (\xone mm + \diameter pt , \ytwo mm)
 %~ arc (270:180:\diameter pt)
	%~ -- (\xone mm, \yone mm);
%~ \end{tikzpicture}
	%~ \begin{textblock}{11.9}[0,0](4,2.73)
         %~ \changefont{phv}{m}{n}
         %~ \raggedleft{ }
     %~ \end{textblock}
	%~ \changefont{phv}{m}{n}	% helvetica	
	%~ \vspace*{2.925cm}
	%~ \begin{center}
		%~ \Huge{\mytitle}
		%~ \vspace*{1cm}\\
		%~ \Large{\myname}\\
		%~ \vspace*{2cm}
		%~ \vspace*{11cm}
		%~ \Large{Diplomarbeit\\
			%~ KIT, \mydepartment
			%~ \\
			%~ \myinstitute}
	%~ \end{center}
%~ \selectlanguage{english}
%~ \end{titlepage}
%\selectlanguage{english}
\begin{titlepage}
%\selectlanguage{USenglish}
% bg shape
\begin{tikzpicture}[overlay]
\draw[color=gray]  
 		 (\xone mm, \yone mm)
  -- (\xtwo mm, \yone mm)
 arc (90:0:\diameter pt) 
  -- (\xtwo mm + \diameter pt , \ytwo mm) 
	-- (\xone mm + \diameter pt , \ytwo mm)
 arc (270:180:\diameter pt)
	-- (\xone mm, \yone mm);
\end{tikzpicture}
	\begin{textblock}{10}[0,0](2,1.2)
		\includegraphics[width=.3\textwidth]{logos/KITLogo_RGB.pdf}
	\end{textblock}
	\begin{textblock}{11.9}[0,0](2,1.5)
         \changefont{phv}{m}{n}
         \raggedleft{\Large{IEKP-KA/2012-15}}
     \end{textblock}
	\changefont{phv}{m}{n}	% helvetica	
	\vspace*{3.225cm}
	\begin{changemargin}{0.5cm}{0cm}
	\begin{center}
		\Huge{mytitle}
		\vspace*{1cm}\\
		\Large{myname}\\
		\vspace*{2cm}
		\Large{
			Master Thesis of
		}\\
		\vspace*{3.5cm}
		\Large{
			mydepartment
			\\
			myinstitute
		}
	\end{center}
	\vspace*{0.6cm}
\Large{
\begin{center}
\begin{tabular}[ht]{l c l}
 Reviewer: & \hfill  & \reviewerone\\
 Second reviewer: & \hfill  & \reviewertwo\\
  %\iflanguage{english}{Advisor}{Betreuender Mitarbeiter}: & \hfill  & \advisor\\
\end{tabular}
\end{center}
}


\vspace{5cm}
\begin{center}
\large{myplace, timeend}
\end{center}
\end{changemargin}

\begin{textblock}{10}[0,0](2,15.45)
\tiny{ 
KIT -- University of the State of Baden-Wuerttemberg and National Laboratory of the Helmholtz Association
}
\end{textblock}

\begin{textblock}{11.9}[0,0](12,15.45)
\large{
	\raggedleft{\textbf{www.kit.edu}}
}
\end{textblock}
\selectlanguage{english}
\end{titlepage}
%\KOMAoptions{twoside=true,DIV=last}

\cleardoublepage
\vspace*{36\baselineskip}
\hbox to \textwidth{\hrulefill}
\par
\iflanguage{english}{I declare that I have developed and written the enclosed thesis completely by myself, and have not used sources or means without declaration in the text.}{Ich versichere wahrheitsgem\"a\ss, die Arbeit selbstst\"andig angefertigt, alle benutzten Hilfsmittel vollst\"andig und genau angegeben und alles kenntlich gemacht zu haben, was aus Arbeiten anderer unver\"andert oder mit Ab\"anderungen entnommen wurde.}

\textbf{PLACE, DATE}
\vspace{1.5cm}

\dotfill\hspace*{8.0cm}\\
\hspace*{2cm}(\textbf{YOUR NAME}) %center name with hspace

\thispagestyle{empty}
\blankpage


%% -------------------
%% |   Directories   |
%% -------------------
\tableofcontents
\blankpage


%% -----------------
%% |   Main part   |
%% -----------------
\mainmatter
\pagenumbering{arabic}
%% introduction.tex
%%

%% ==============================
\chapter{Introduction}
\label{ch:Introduction}
%% ==============================

\cite{becker2008a}
\dots


%% content.tex
%%

%% ==============
\chapter{Content Chapter 1}
\label{ch:Content1}
%% ==============

The content chapters of your thesis should of course be renamed. How many chapters you need to write depends on your thesis and cannot be said in general. 

Check our the examples theses in the Wiki. 

Of course, you can split this .tex file into several files if you prefer. 


%% ===========================
\section{Section 1}
\label{ch:Content1:sec:Section1}
%% ===========================

\dots


%% ===========================
\section{Section 2}
\label{ch:Content1:sec:Section2}
%% ===========================

\dots


%% content.tex
%%

%% ==============
\chapter{Content Chapter 2}
\label{ch:Content1}
%% ==============

\dots


%% ===========================
\section{Section 1}
\label{ch:Content2:sec:Section1}
%% ===========================

\dots


%% ===========================
\section{Section 2}
\label{ch:Content2:sec:Section2}
%% ===========================

\dots

Add additional content chapters if required. 
%% evaluation.tex
%%

%% ==================
\chapter{Evaluation}
\label{ch:Evaluation}
%% ==================

\dots


%% ===============================
\section{Section 1}
\label{ch:Evaluation:sec:Section1}
%% ===============================

\dots


%% ===============================
\section{Section 2}
\label{ch:Evaluation:sec:Section2}
%% ===============================

\dots


%% ===============================
\section{Section 3}
\label{ch:Evaluation:sec:Section3}
%% ===============================

\dots
%% conclusion.tex
%%

%% ==================
\chapter{Conclusion}
\label{ch:Conclusion}
%% ==================

\dots



%% --------------------
%% |   Bibliography   |
%% --------------------
\cleardoublepage
\phantomsection
\addcontentsline{toc}{chapter}{\bibname}

\iflanguage{english}
{\bibliographystyle{IEEEtranSA}}	% english style
{\bibliographystyle{babalpha-fl}}	% german style
												  
% Use IEEEtran for numeric references
%\bibliographystyle{IEEEtranSA})

\bibliography{thesis}


%% ----------------
%% |   Appendix   |
%% ----------------
\cleardoublepage

%% appendix.tex
%%

%% ==============================
%\chapter{Appendix}
%\label{ch:Appendix}
%% ==============================

\appendix

\iflanguage{english}
{\addchap{Appendix}}	% english style
{\addchap{Anhang}}	% german style


\section{First Appendix Section}
		\label{Anhang-Implementierung}
		
\setcounter{figure}{0}
		
\begin{figure} [ht]
  \centering
   ein Bild
  \caption{A figure}
  \label{fig:BPMNBeispiela}
\end{figure}


\dots






\end{document}
